\section*{Sets and Indices}

\[
T=\{0,1,2,3,4,5\}
\]
monthly periods, corresponding respectively to January, February, March, April, May, and June.

\[
T^{c}=\{3,4,5\}
\]
contract months (April--June).

\section*{Parameters}

For each \(t\in T\), let \(d_t\) denote forecast demand in month \(t\) (code: \texttt{demand[t]}), with
\[
(d_0,d_1,d_2,d_3,d_4,d_5)=(20000,40000,42000,35000,19000,18500).
\]

Scalar parameters from \texttt{general\_parameters.csv}:

\[
W^0=1000
\]
initial workforce (code: \texttt{W0}),

\[
I^0=15000
\]
initial inventory (code: \texttt{I0}),

\[
p=300
\]
sales price per unit (code: \texttt{sp}),

\[
c^{rm}=90
\]
raw material cost per in-house unit (code: \texttt{rmc}),

\[
c^{out}=200
\]
outsourcing cost per unit (code: \texttt{oc}),

\[
c^{inv}=15
\]
inventory holding cost per unit-month (code: \texttt{hc}),

\[
c^{bo}=35
\]
backorder cost per unit-month (code: \texttt{bc}),

\[
\ell=5
\]
labor hours required per in-house unit (code: \texttt{lr}),

\[
h^{reg}=160
\]
regular labor hours per worker per month (code: \texttt{rh}),

\[
w^{reg}=30
\]
regular wage per labor hour (code: \texttt{rw}),

\[
h^{ot}=20
\]
maximum overtime hours per worker per month (code: \texttt{otl}),

\[
w^{ot}=40
\]
overtime wage per hour (code: \texttt{ow}),

\[
c^{hire}=5000
\]
hiring cost per worker (code: \texttt{hire\_cost}),

\[
c^{fire}=8000
\]
firing cost per worker (code: \texttt{fire\_cost}),

\[
I^{\min}_{6}=10000
\]
minimum ending inventory in June (code: \texttt{term\_inv}),

\[
B^{\max}_{6}=0
\]
maximum ending backorders in June (code: \texttt{term\_bo}),

\[
Q^{\min}=15000
\]
minimum contract-qualified units in each contract month (code: \texttt{contract\_min\_units}),

\[
M^{c}=60000
\]
big-\(M\) constant for contract linkage (code: \texttt{contract\_bigM}).

Also define total horizon demand
\[
D^{\mathrm{tot}}=\sum_{t\in T} d_t.
\]

\section*{Decision Variables}

For each \(t\in T\):

\[
W_t \ge 0
\]
workforce in month \(t\) (code: \texttt{W[t]}),

\[
H_t \ge 0
\]
workers hired in month \(t\) (code: \texttt{H[t]}),

\[
F_t \ge 0
\]
workers fired in month \(t\) (code: \texttt{F[t]}),

\[
P_t \ge 0
\]
total in-house production (code: \texttt{P[t]}),

\[
P^{reg}_t \ge 0
\]
regular-time in-house production (code: \texttt{P\_reg[t]}),

\[
P^{ot}_t \ge 0
\]
overtime in-house production (code: \texttt{P\_ot[t]}),

\[
OT_t \ge 0
\]
total overtime hours used (code: \texttt{OT\_total[t]}),

\[
O_t \ge 0
\]
outsourced units (code: \texttt{O[t]}),

\[
I_t \ge 0
\]
ending inventory (code: \texttt{Inv[t]}),

\[
B_t \ge 0
\]
ending backorders (code: \texttt{B[t]}),

\[
C_t \ge 0
\]
contract-qualified units (code: \texttt{C[t]}),

\[
z_t \in \{0,1\}
\]
contract activation indicator (code: \texttt{z[t]}).

All nonbinary variables are continuous, matching the implementation.

\section*{Objective}

The code maximizes net profit:
\[
\max \; p\,D^{\mathrm{tot}}
-\sum_{t\in T}\Big(
c^{rm} P_t
+c^{out} O_t
+c^{inv} I_t
+c^{bo} B_t
+w^{reg} h^{reg} W_t
+w^{ot} OT_t
+c^{hire} H_t
+c^{fire} F_t
\Big).
\]

Because the implementation sets revenue equal to \(p\) times total forecast demand, revenue is a constant term.

\section*{Constraints}

\paragraph{Workforce balance}
For \(t=0\):
\[
W_0 = W^0 + H_0 - F_0.
\]
For \(t=1,\dots,5\):
\[
W_t = W_{t-1} + H_t - F_t.
\]

\paragraph{Regular-time capacity}
For all \(t\in T\):
\[
\ell\, P^{reg}_t \le h^{reg} W_t.
\]

\paragraph{Overtime production-hour linkage}
For all \(t\in T\):
\[
\ell\, P^{ot}_t = OT_t.
\]

\paragraph{Overtime capacity}
For all \(t\in T\):
\[
OT_t \le h^{ot} W_t.
\]

\paragraph{Production decomposition}
For all \(t\in T\):
\[
P_t = P^{reg}_t + P^{ot}_t.
\]

\paragraph{Inventory-backorder balance}
For \(t=0\):
\[
I_0 - B_0 = I^0 + P_0 + O_0 - d_0.
\]
For \(t=1,\dots,5\):
\[
I_t - B_t = I_{t-1} - B_{t-1} + P_t + O_t - d_t.
\]

\paragraph{Contract accounting for April--June}
For each \(t\in T^c\):
\[
C_t \le P^{reg}_t,
\]
\[
C_t \le d_t + B_{t-1},
\]
\[
C_t \ge Q^{\min},
\]
\[
C_t \le M^{c} z_t.
\]

\paragraph{No contract activity before April}
For \(t=0,1,2\):
\[
C_t = 0,
\]
\[
z_t = 0.
\]

\paragraph{Terminal conditions}
\[
I_5 \ge I^{\min}_{6},
\]
\[
B_5 \le B^{\max}_{6}.
\]

\section*{Notes on Implementation}

The formulation above matches the implemented mixed-integer linear program in \texttt{current\_heuristic.py}, solved by Gurobi through the PuLP-compatible interface.

The contract logic is implemented through the auxiliary variable \(C_t\) and binary \(z_t\) for April--June. Since \(C_t \ge Q^{\min}=15000>0\) and \(C_t \le M^c z_t\), feasibility forces \(z_t=1\) in each contract month; the binary variable is therefore redundant in effect, but it is present in the code and is included here exactly as implemented.

The code interprets the monthly contract cap based on
\[
d_t + B_{t-1},
\]
i.e., current-month demand plus prior-month backlog. It does not introduce a separate shipment-allocation variable tying \(C_t\) to actual units delivered from regular production; instead, \(C_t\) is only upper-bounded by regular production and by this ``net demand'' expression.

All workforce, hiring, and firing variables are continuous in the implementation rather than integer-valued, so the model permits fractional workers. Revenue is treated as the constant \(p\sum_t d_t\), relying on the terminal condition \(B_5=0\) to ensure all demand is eventually satisfied within the horizon.
